Thirteen further simulations mapped the variables that Stages~2 and 4 had identified as controlling (Figs.~\ref{fig:deep_hierarchy}--\ref{fig:deep_scaling}).

\paragraph{Vein width, scale ratio and inter-level interactions.} At matched porosity the strength of the two-level mesh rises monotonically with the width of its veins, 13.2/12.8/13.7/14.8/15.6\,N/m for $W=4/8/12/16/20$\,\AA, toward the coarse-only limit (17.0\,N/m), and the mode turns from progressive to abrupt at $W=20$\,\AA{} when the fine pores have shrunk to a few atoms; work to failure stays at 1.3--1.5\,J/m$^2$ throughout. The scale ratio at fixed vein geometry has no effect (fine period 8 vs 12\,\AA: 12.7 vs 12.8\,N/m), and a larger domain with wider veins (60/12\,\AA) gives 13.0\,N/m. Disordering the fine level costs the nested mesh the same 2--3\,N/m it costs a single-scale mesh (12.8 $\to$ 10.3). The single favourable interaction is hierarchy $\times$ anisotropy: elongating the fine pores along the load inside the nested mesh gives 17.5\,N/m, $W=2.0$\,J/m$^2$ and progressive failure with 0.26 failure strain -- the only nested design that exceeds the coarse-only mesh, because it combines straight fine ligaments (net-section strength close to the intrinsic value) with many parallel load paths (sequential failure).

\paragraph{Disorder.} The Voronoi regularity sweep (0/0.3/0.6/0.8/1.0: 9.2$\pm$0.5, 10.5, 11.2$\pm$2.8, 6.6, 13.3\,N/m) has no monotonic trend once the seed scatter is taken into account; only the fully regular (hexagonal-seed) network is reproducibly strong. The random networks fail from their thinnest ligament (localisation $L\approx0.4$--0.6), which is why their net-section strength (12--20\,N/m) lies below that of ordered meshes of the same ligament width (Fig.~\ref{fig:deep_scaling}, right).

\paragraph{Porosity scaling and the net-section rule.} For load-parallel slits the strength is independent of porosity (25.6/25.4/25.2\,N/m for $\phi=0.10/0.20/0.22$; the slit family cannot exceed $\phi\approx0.22$ at 4\,\AA{} slit width, so the two nominal 0.3/0.4 members are replicates of the 0.22 design and reproduced it exactly), because lengthening the slits does not change the net section perpendicular to the load. The square coarse mesh follows $\sigma\approx31\times\mathrm{minSF}$ (21.8/17.0/14.2\,N/m at $\phi=0.1/0.2/0.3$), whereas the fine round-pore mesh follows the lower line $\sigma\approx17\times\mathrm{minSF}$ (19.6/11.7/9.4/7.3\,N/m at 0.1--0.4) with its net-section strength itself decreasing as the ligaments thin. Across all porous designs the net-section strength is 45--55\% of the intrinsic strength for ligaments of 10--13\,\AA, 70--85\% for 25--30\,\AA, with load-parallel slits at the top and random networks scattered below (master plot).
